\section*{Chapter \ref{ch:rs:the-research-process} --- The Research Process}

\begin{solution}{exo:rs:the-research-process:1}
$\mathrm{PPV} = 0.1 \times 0.8/(0.1 \times 0.8 + 0.9 \times 0.05) = 0.08/0.125 = 0.64$. With a 2\% prior, $0.016/(0.016 + 0.049) = 0.246$:
three passes in four are false.
\end{solution}

\begin{solution}{exo:rs:the-research-process:2}
$1 - \Phi(2)^{20} = 1 - 0.97725^{20} = 0.369$; on average $20 \times 0.02275 = 0.455$ pass.
\end{solution}

\begin{solution}{exo:rs:the-research-process:3}
$\Phi\bigl((-0.5 - 1.2)\sqrt{0.5}\bigr) = \Phi(-1.20) = 0.115$: about one half-year in nine would be this bad. Not decisive on its own.
\end{solution}

\begin{solution}{exo:rs:the-research-process:4}
$Y = ((1.645 + 1.282)/0.8)^2 = 13.4$ years, against 2.7 years for 1.5 at 80\% power: the time grows with the inverse square of
the Sharpe ratio.
\end{solution}

\begin{solution}{exo:rs:the-research-process:5}
$30 + \binom{30}{2} = 30 + 435 = 465$ tests. Bonferroni needs $p \le 0.05/465 = 1.08 \times 10^{-4}$, that is $t \ge 3.70$.
\end{solution}

\begin{solution}{exo:rs:the-research-process:6}
Six real and 144 false ideas. After stage one, 3.6 real and 7.2 false (PPV 0.333); after stage two, 2.16 and 0.36 (0.857);
after stage three, 1.30 and 0.018 (0.986). About 1.3 real strategies and one false one every 55 years reach production.
\end{solution}

\begin{solution}{exo:rs:the-research-process:7}
\texttt{p\_best(2, 1000, 2, rho)} gives 0.420, 0.167 and 0.030 for $\rho = 0.3$, 0.6 and 0.9. Twenty independent variants give 0.046;
a thousand variants equicorrelated at $\rho = 0.85$ give the same. Even at 0.9 a thousand variants are worth about thirteen
independent ones: correlation reduces a search but rarely to one trial.
\end{solution}

\begin{solution}{exo:rs:the-research-process:8}
The holdout stopped being a holdout when it was used to tune: the 2.1 is an in-sample number, and the only independent evidence
is the 1.4 measured before tuning, itself one look at a holdout that was then reused. Two parameters tuned on three years are also
trials the log must count. The honest status is one exploratory stage passed; a new, untouched period (or live shadow trading)
is needed before the claim ``two independent stages'' is true.
\end{solution}

\begin{solution}{pb:rs:the-research-process:1}
\textbf{1.} Approximately $\mathcal N(0, 1/2)$: standard deviation 0.707. \textbf{2.} $1 - \Phi(2\sqrt2) = 1 - \Phi(2.83) = 0.00234$.
\textbf{3.} $1\,000 \times 0.00234 = 2.34$. \textbf{4.} The question is the probability that at least one impresses (and which one is
shown); correlation changes that probability, not the expectation. \textbf{5.} $1 - (1 - 0.00234)^{20} = 0.0458$.
\textbf{6.} By \cref{prop:rs:the-research-process:best} with $\rho = 0.6$: 0.0262. \textbf{7.} $1 - (1 - 0.0262)^{50} = 0.735$.
\textbf{8.} $1 - \Phi(2.83)^{1\,000} = 0.904$. \textbf{9.} 2.30, above the impressive threshold.
\textbf{10.} $\Phi\bigl(2/\sqrt{252} \times \sqrt{503}\bigr) = 0.998$. \textbf{11.} With the benchmark the expected best of twenty null
ratios: 0.822. \textbf{12.} With a thousand: 0.336, a coin flip that the strategy is better than the best of the search.
\textbf{13.} The count must cover every trial that could have been shown; correlation reduces the effective number, and the
deflated ratio with the raw count is then conservative, but the effective count needs the variants' correlations, which only
the log can provide. \textbf{14.} $2\sqrt Y \ge z_{0.99999} = 4.265$, so $Y \ge 4.55$ years.
\textbf{15.} $1 - \Phi(1) = 0.159$. \textbf{16.} Power $\Phi(1.5 - 1) = 0.691$; $\mathrm{PPV} = 0.05 \times 0.691/(0.05 \times 0.691 + 0.95
\times 0.159) = 0.187$: a one-year holdout at this threshold is a weak stage. \textbf{17.} $1 - (1 - 0.159)^4 = 0.499$: a coin
flip. \textbf{18.} \emph{Named result:} a worthless search of twenty correlated variants a week produces at least one backtest
with a Sharpe ratio of 2 within a year with probability 0.735, and that backtest's deflated Sharpe ratio with the logged
thousand trials is 0.336. \textbf{19.} The hypothesis and the date it was registered; the \omterm{def:rs:the-research-process:log}{trial count} of the family and the
distribution of all trials' results; the evidence beyond the in-sample fit (later data, side predictions, a holdout).
\textbf{20.} It records the search, and the search is the largest unmodelled risk in a backtest.
\end{solution}

\begin{solution}{iq:rs:the-research-process:1}
How many variants were tried, and how correlated were they (the \omterm{def:rs:the-research-process:log}{trial count}, and so the deflated ratio)? What is the \omterm{def:rs:the-research-process:hypothesis}{economic
rationale}, and was it written before the test? What evidence exists beyond the in-sample fit (later data, other markets, side
predictions), and are the costs, borrow and capacity realistic?
\iqlookfor{the search before the statistics; a rationale with testable side predictions; costs.}
\end{solution}

\begin{solution}{iq:rs:the-research-process:2}
By Bayes, the share of true strategies among those that pass is $\pi(1 - \beta)/(\pi(1 - \beta) + (1 - \pi)\alpha)$. With a 5\% prior,
5\% size and 50\% power it is 0.345: most passes are false whatever the size, unless the prior rises or independent stages
are added. Size controls the false positives per false idea tested, not the proportion of false strategies among the winners.
\iqlookfor{the positive predictive value, and that priors and stages dominate.}
\end{solution}

\begin{solution}{iq:rs:the-research-process:3}
Any use of it that influences a choice: tuning on it, looking at it and going back to change the model, re-running it after a
disappointing result, or choosing the sample period with knowledge of it. Each look raises the holdout's effective size (five
looks at 5\% give 22.6\%).
\iqlookfor{that a holdout is spent by influence, not by being read once.}
\end{solution}

\begin{solution}{iq:rs:the-research-process:4}
Not on that evidence alone: if the true ratio were 1.5, a year at zero or below has probability $\Phi(-1.5) = 0.067$. Look at what
is measured faster than returns: are fills, costs and turnover what the simulator said? Is the signal's live information
coefficient in line with the backtest? Apply the \omterm{def:rs:the-research-process:gate}{kill criterion} fixed before going live; if none was fixed, fix one now and
include a minimum time.
\iqlookfor{the noise in a one-year Sharpe ratio, and diagnostics beyond P\&L.}
\end{solution}

\begin{solution}{iq:rs:the-research-process:5}
Per entry: family, kind (hypothesis or trial), caller-supplied timestamp, parameters, metrics, code version, data snapshot and
the digest of the previous entry, with the entry's own digest a hash of all of these. Append-only storage; verification walks
the chain. One hypothesis per family registered before any trial; trials before it flag the family post hoc.
\iqlookfor{hash chaining, reproducible timestamps, and counting trials per hypothesis family.}
\end{solution}

\begin{solution}{iq:rs:the-research-process:6}
Sellers or buyers demanding immediacy push prices away from value, and liquidity providers who absorb the flow are paid by the
subsequent reversal. Predictions: the reversal is larger in less liquid stocks; larger when market-wide liquidity is scarce
(high volatility, dealer balance sheets constrained); smaller or absent when the move carries news (earnings days), and
shrinking as the market's cost of immediacy falls over the years.
\iqlookfor{a rationale that says who pays, and side predictions that could refute it.}
\end{solution}
